\section{Conclusions}

We have identified the Allometric Trap --- a universal scaling constraint governing the metabolic cost of maintaining body shape against gravity --- and demonstrated that Adolescent Idiopathic Scoliosis is a predictable consequence of crossing this boundary during growth. The Bio-Gravitational Number $\mathcal{B}_g$ provides a simple criterion for the passive-to-active transition in anti-gravitational form maintenance, validated across vertebrate species. Exploratory AlphaFold structural analysis generates hypotheses regarding the molecular architecture of this vulnerability, suggesting a potential demand-supply structural gap between mechanosensory and metabolic proteins. The Delay Differential Equation (DDE) control framework unifies mathematical control theory, biomechanics, and developmental metabolism under a single structure, generating three falsifiable predictions testable with existing technology. These results suggest that scoliosis prevention strategies should shift from purely mechanical approaches to metabolic interventions targeting the Energy Deficit Window --- closing the gap between what the spine costs and what the body can afford.
